\section{Insights and Conclusions}

\subsection{Insights}

\subsubsection*{Tuition alone explains little}

Model 1 establishes that the unconditional relationship is real but thin. The pooled tuition elasticity of 0.121 (95\% CI [0.113, 0.129], p $<$ .001) is statistically unambiguous across 2,894 institutions, yet it accounts for only 18.7\% of the variation in log earnings. Four-fifths of the difference in what graduates earn has nothing to do with what their institution charges. The Spearman coefficient of 0.464 from the preliminary inference tells the same story in rank terms: a moderate, reliably non-zero association, not a deterministic one.

\subsubsection*{Sector is the stratifying variable, not the price}

Adding institutional control and its interaction with tuition raises the adjusted R-squared from 0.187 to 0.397, and the joint HC3 Wald test on the interaction terms is decisive (statistic 30.122, p $<$ .001). Sector alone contributes roughly the same explanatory weight as tuition itself. The unadjusted slopes appear to invert the intuitive ordering: for-profits post the steepest elasticity at 0.400, ahead of private nonprofits at 0.277 and public institutions at 0.194. Taken at face value, this would suggest that the sector with the weakest earnings outcomes converts tuition into earnings most efficiently.

\subsubsection*{That inversion is composition, not return}

Model 3 dissolves it. Once undergraduate enrollment, predominant degree, and state fixed effects are held constant, the for-profit elasticity falls from 0.400 to 0.119, meaning roughly 70\% of the unadjusted coefficient was absorbed by observable institutional structure. For-profits cluster in particular states, enroll small student bodies, and award predominantly sub-baccalaureate credentials; each of these depresses earnings independently of price, and the unadjusted interaction was reading those features as a tuition effect. After adjustment the ordering reverses to 0.189 for public, 0.147 for private nonprofit, and 0.119 for for-profit institutions. The sector difference survives, but only marginally: the joint interaction test returns p = .0465, and the restricted selectivity sample cannot reproduce it at all (p = .155). \textcite[pp.~35--37]{muse2024} document the same admission-coverage limitation for institution-level earnings regressions.

\subsubsection*{Where an institution sits matters more than what it charges}

The control block adds 18.5 percentage points of adjusted explanatory power beyond Model 2, almost exactly what tuition contributed on its own, and the controls are jointly significant at p $<$ 10$^{-265}$. For a prospective student, the practical reading is that geography and credential type carry at least as much information about likely earnings as the sticker price does. Even so, 41.8\% of the variation remains unexplained (adjusted R-squared of 0.582), which is the expected ceiling for institution-level aggregates that cannot observe student background, field of study, or local labour demand. \textcite[pp.~45--47]{faber2017} identifies the same causal and omitted-variable boundaries for institution-level College Scorecard regressions.

\subsubsection*{Elasticity ranking and dollar return are not the same ranking}

Elasticities are scale-free, so a larger coefficient does not mean a better deal. Priced at each sector's median institution, a 10\% tuition increase costs \$598 at a public institution and returns roughly \$822 in predicted annual earnings, about \$1,374 per \$1,000 of additional tuition. The same proportional increase costs \$3,393 at a private nonprofit and returns roughly \$767, about \$226 per \$1,000, and \$1,633 at a for-profit for roughly \$447, about \$274 per \$1,000. Public institutions convert price into predicted earnings around six times more efficiently than private nonprofits, precisely because they start from a low tuition base. This is a descriptive per-dollar ratio rather than a payback calculation, since tuition recurs annually across a degree while the earnings figure is a single-year median at ten years post-entry.

\subsubsection*{The sector gap is a level gap, not a slope gap}

The median private nonprofit graduate out-earns the median for-profit graduate by \$14,718. Closing that gap through tuition alone, at the for-profit sector's own adjusted elasticity of 0.119, would require the median for-profit institution to charge roughly sixteen times its current tuition, above \$260,000. The sectors sit on different intercepts, and no plausible movement along a within-sector slope crosses the distance between them. Combined with the debt-to-earnings ratios of 0.29 for public graduates against 0.45 and 0.43 for the other two sectors, the net return picture favours the public sector on both the numerator and the denominator.

\subsection{Conclusion}

This study sought to determine whether a significant relationship exists between tuition and post-graduation earnings, and the results provide a clear answer. There is a statistically significant, moderate positive relationship between in-state tuition and median earnings, supported by a Spearman correlation coefficient of $r$ = 0.4641 and a p-value of 1.53 $\times$ 10$^{-154}$, which is far below the 0.05 significance level. This leads to the rejection of the null hypothesis and confirms that higher tuition is, on average, associated with higher earnings.

However, this relationship is not consistent across all institutional sectors. Public institutions exhibit the strongest and most reliable relationship, indicating a more predictable return on investment. In contrast, for-profit institutions show the weakest relationship, suggesting that higher tuition does not consistently lead to higher earnings in this sector. Private nonprofit institutions, while associated with higher tuition and higher average earnings, display greater variability, meaning that outcomes are less predictable despite higher costs.

These findings carry important implications. For students, the results highlight that tuition alone should not be used as a proxy for future earnings; instead, the type of institution and the consistency of outcomes must be considered. For policymakers, the variation across sectors emphasizes the need for improved transparency and accountability, particularly in sectors where the link between cost and return is weaker. Providing accessible and reliable data can help ensure that educational investments lead to meaningful economic outcomes.

Overall, the study confirms that a significant relationship exists, that it differs by sector, and that it has real implications for both individual decision-making and higher education policy.
